\exnumber{01}
\extitle{UX V - Ex 2 - Plan de site}
\exfiles{ex01/fichier Figma et PDF, nomme Name\_FirstName\_DeliverableName\_Date\_VersionNumber}
\exauthorize{Ceux nommés dans l'exercice}
\makeheaderfiles

\textbf{Instructions :}

Restez sur le site web que vous avez choisi dans l'Ex.1 et concevez un plan de site. Le plan de site doit inclure toutes les pages du site ainsi que les connexions et la hiérarchie.

\textbf{Liste de sites :}

\begin{itemize}
\item Universal Music Latin America \url{https://www.universalmusica.com/}
\item Delta Records \url{https://deltarecords.net/}
\item Zoho Music \url{https://www.zohomusic.com/}
\item Luaka Bop \url{https://www.luakabop.com/}
\item Rimas Music \url{https://rimasmusic.com/}
\item D'autres... tant que c'est un label de musique sud-américain ;-)
\end{itemize}

\textbf{Livrables :}

\begin{itemize}
\item Une page PDF du plan de site.
\item Déposez le PDF sur Github, nommé ainsi : « Name\_FirstName\_DeliverableName\_Date\_VersionNumber ». Exemple : « Doe\_John\_MidFiPrototype\_05242024\_V1 ».
\item Ajoutez un titre dans le document.
\item Soignez l'aspect global du document. Il doit être clair, simple et lisible. Vous pouvez regarder canva.com, sans surcharger le dessin de détails.
\end{itemize}

\textbf{Pistes :}

\begin{itemize}
\item Commencez par dessiner sur papier, puis passez à un outil numérique.
\item Veuillez intituler votre document ! Titre et nom du site web.
\end{itemize}

\textbf{Ressources :}

\begin{itemize}
\item Comment créer un bon plan de site \url{https://www.smokeylemon.com/blog/how-to-create-a-good-sitemap/}
\end{itemize}

\textbf{Outils recommandés :}

\begin{itemize}
\item Figma
\item Adobe Illustrator
\item etc.
\end{itemize}

\textbf{Citation :}

« Les systèmes d'organisation, de recherche et de navigation qui aident les gens à accomplir des tâches, trouver ce dont ils ont besoin, et comprendre ce qu'ils ont trouvé » Peter Morville | Information Architecture for the WWW

